\section{Related Work}

Prior work relevant to this paper spans four overlapping areas: classical retrieval-augmented generation, agentic RAG, multi-agent large language model systems, and educational AI systems for tutoring, feedback, and adaptive support \cite{chevalier2024sciencetutors,liu2026agenttutor,macina2025mathtutorbench,srinivasa2025tutorbench,xu2025edubench}. Classical RAG remains a strong baseline for knowledge-grounded tasks because it is simple, interpretable, and often effective when answers depend on external evidence. At the same time, more recent agentic systems suggest that retrieval alone may be insufficient for tasks that require iterative planning, critique, or tool use \cite{nguyen2025marag,radeva2025coordination}. Multi-agent approaches further extend this by introducing role specialization and coordination across multiple components, rather than relying on a single retrieve-then-generate pipeline \cite{liu2023agentbench,nguyen2025marag,radeva2025coordination}.

In educational settings, this distinction matters because many tasks are not defined only by factual correctness. Tutoring dialog, rubric-based feedback, misconception diagnosis, and study planning also require pedagogical adaptation, instructional sequencing, and sensitivity to the state of the learner \cite{chevalier2024sciencetutors,liu2026agenttutor,macina2025mathtutorbench,srinivasa2025tutorbench,xu2025edubench}. Education-facing benchmarks such as EduBench, TutorEval, ScienceQA, MathTutorBench, and TutorBench emphasize pedagogical quality, adaptivity, explanation, and instructional alignment rather than only factual correctness \cite{xu2025edubench,chevalier2024sciencetutors,lu2022scienceqa,macina2025mathtutorbench,srinivasa2025tutorbench}. By contrast, transferable benchmarks such as HotpotQA, SCROLLS, LongBench v2, FEVER, Wizard of Wikipedia, BEIR, and AgentBench isolate mechanism-level abilities such as retrieval depth, evidence justification, long-context handling, grounded dialogue, retrieval robustness, fact verification, and workflow coordination \cite{yang2018hotpotqa,shaham2022scrolls,bai2024longbench2,thorne2018fever,dinan2019wizard,thakur2021beir,liu2023agentbench}.  This motivates the paper's central distinction between retrieval-heavy and coordination-heavy task regimes, and frames the comparison among classical RAG, agentic RAG, and multi-agent systems as a question of task fit rather than universal superiority.
